\documentclass[10pt,a4paper]{article}
\usepackage[utf8]{inputenc}
\usepackage[T1]{fontenc}
\usepackage[spanish]{babel}
\usepackage[margin=20mm]{geometry}
\usepackage{amsmath,amssymb,booktabs,hyperref}
\usepackage{listings}
\lstset{basicstyle=\ttfamily\small,breaklines=true,frame=single}
\title{Material viscoelástico SLS en MUSCLE\\\large Integración por punto de Gauss con un shell T3}
\author{Proyecto XFEM / ContactProj}
\date{26 de octubre de 2026}
\begin{document}
\sloppy
\maketitle
\section*{Modelo y ubicación}
Se incorpora \texttt{Viscoelastic\_sls} en
\texttt{src/viscoelasticity/muscle\_viscoelastic\_sls.f90}, compilado como parte de
\texttt{libmuscle}. Es una implementación nueva para este proyecto, no un modelo
preexistente de la biblioteca. Hereda de \texttt{Elasticity\_linear} y utiliza los
tensores y la interfaz transitoria de MUSCLE. No cambia los materiales existentes.

El sólido lineal estándar consta de un resorte de equilibrio en paralelo con una
rama Maxwell (resorte y amortiguador en serie). En pequeñas deformaciones:
\begin{align}
 \boldsymbol\sigma &= \mathbb C_\infty:\boldsymbol\varepsilon+\mathbf q,\\
 \dot{\mathbf q}+\frac{\mathbf q}{\tau}
 &= r\,\mathbb C_\infty:\dot{\boldsymbol\varepsilon}.
\end{align}
$\mathbb C_\infty$ se obtiene con el material elástico de MUSCLE;
$r=E_M/E_\infty$ y $\tau$ es el tiempo de relajación.
Ambas ramas tienen el mismo coeficiente de Poisson: se relajan las respuestas
volumétrica y desviadora con el mismo tiempo. No es un modelo de plasticidad.

\section*{Actualización local y tangente}
Para deformación lineal durante un incremento $\Delta t$:
\begin{align}
 a&=e^{-\Delta t/\tau}, & b&=\frac{1-a}{\Delta t/\tau},\\
 \mathbf q_{n+1}^{\mathrm{trial}}&=a\mathbf q_n+
 r b\,\mathbb C_\infty:(\boldsymbol\varepsilon_{n+1}-\boldsymbol\varepsilon_n),\\
 \boldsymbol\sigma_{n+1}&=\mathbb C_\infty:\boldsymbol\varepsilon_{n+1}
 +\mathbf q_{n+1}^{\mathrm{trial}}, &
 \mathbb C_{\mathrm{alg}}&=(1+rb)\mathbb C_\infty.
\end{align}
Se usa una serie de Taylor para evaluar $b$ cuando $\Delta t/\tau$ es pequeño.
Cada punto de Gauss tiene su propio \texttt{Sls\_internal}, que almacena los seis
componentes de $\mathbf q$. El FEM conserva también la deformación aceptada para
construir el incremento. La interfaz Fortran usa cortantes tensoriales; el puente C
recibe cortantes ingenieriles y realiza la conversión.

\section*{Uso desde Fortran}
\begin{lstlisting}[language=Fortran]
use muscle_viscoelastic_sls
! eps, deps, sigma: ten_3D2Osym
! C: ten_3D4O3sym
! Declarar un estado independiente por punto de Gauss:
type(Viscoelastic_sls) :: mat
type(Sls_internal) :: history
call mat%set_parameters(1.0d7, 0.3d0)
mat%ratio = 1.0d0
mat%tau = 0.5d0
sigma = mat%stress(eps, deps, dt, history)
C = mat%dstress_dstrain(eps, deps, dt, history)
! Solo despues de converger:
call mat%stress_and_update(sigma, eps, deps, dt, history)
\end{lstlisting}
\newpage
\section*{Integración en el shell}
ContactProj proporciona las deformaciones generalizadas del T3. Se usa un punto
centroidal en el plano (deformaciones constantes en este elemento lineal) y dos
puntos de Gauss en el espesor:
\[
 z_g=\pm\frac{h}{2\sqrt3},\qquad w_g=\frac h2,\qquad
 \boldsymbol\varepsilon_p(z_g)=\boldsymbol\varepsilon_m+z_g\boldsymbol\kappa.
\]
En cada punto se llama a la rutina Fortran con su historia independiente. El puente
resuelve $\sigma_{zz}=0$, incluyendo la contribución hereditaria, y condensa la
tangente:
\[
 \mathbb C_{ps}=\mathbb C_{pp}^{alg}
 -\mathbb C_{pz}^{alg}(C_{zz}^{alg})^{-1}\mathbb C_{zp}^{alg}.
\]
Luego se integran $\mathbf N=\sum_g w_g\boldsymbol\sigma_p(z_g)$ y
$\mathbf M=\sum_g w_g z_g\boldsymbol\sigma_p(z_g)$. Para cortante transversal se
aplica $\sqrt{5/6}$ a la deformación y a la transformación de tensión; así la
rigidez efectiva lleva el factor $5/6$ y se preserva el trabajo virtual.

El puente actualiza una copia local de la historia. El solver acepta todos los
estados únicamente después de resolver el equilibrio. Las evaluaciones repetidas
durante iteraciones no avanzan la historia. Los resultantes almacenados para VTK
son datos derivados; la memoria constitutiva está en los puntos de Gauss.

\section*{Ensayo de comprobación}
\begin{center}
\begin{tabular}{ll}\toprule
Parámetro & Valor\\\midrule
Dimensiones & $100\times20\times1$ mm\\
$E_\infty$, $E_M$, $\nu$ & 10 MPa, 10 MPa, 0,3\\
$\tau$ & 0,5 s\\
Carga & Rampa a 1\% en 1 s; mantenimiento hasta 4 s\\
Malla / paso & 16 triángulos / 0,05 s\\\bottomrule
\end{tabular}
\end{center}
La fuerza de referencia sube hasta $2{,}86466$ N y se relaja hasta $2{,}00214$ N
al final; su límite es 2 N. Se verifican tracción, malla refinada, tangente,
relajación, descarga, límite elástico y llamadas polimórficas a MUSCLE. La prueba
Fortran comprueba además que evaluar un trial no modifica la historia.

\section*{Alcance y reproducción}
Implementación 3D, CPU y pequeñas deformaciones. La llamada sin tiempo devuelve
la contribución elástica de equilibrio; para el SLS completo se debe usar la
sobrecarga transitoria. Las sobrecargas transitorias 2D se rechazan explícitamente:
el shell usa la rutina 3D con reducción de tensión plana en el puente.
No se afirma validación de flexión libre, contacto, grandes deformaciones ni GPU.
Los parámetros son ilustrativos.

En XFEM: \texttt{./scripts/run.sh}. La secuencia de ParaView se escribe en
\texttt{results/sls\_fortran/vtk/tension.pvd}. El test nativo está en
\texttt{tests/viscoelasticity/test\_viscoelastic\_sls.f90} dentro de MUSCLE.
La rutina usa la API de \texttt{muscle\_elasticity\_base} y
\texttt{muscle\_elasticity\_linear} del mismo repositorio.
\end{document}
